\section{楕円成分、scheme分解と量子係数の符号}
以下さらに$h^6\ne-27$とし、$Y$は滑らかな整曲線とする。
\[
 E_h=V(x^3+y^3+z^3+h^2xyz)\subset\PP^2=V(u,v),\qquad o=[1:-1:0]
\]
と置く。右点加群の点列が$(p_n,p_{n+1})$上の双線形関係を満たす規約を採る。

\begin{theorem}[全点スキームとshift]
安定化した点スキームは、schemeとして
\[
 \operatorname{Point}(B_{h,q})=E_h\sqcup Y,\qquad
 \operatorname{Point}(A_{h,q,f})=E_h\sqcup Z_f
\]
である。shiftは楕円成分で$[1:-1:-h]$によるtranslation $\sigma_h$、
定常成分では恒等となる。
\end{theorem}
\begin{proof}
境界の二次関係は$S(1,-1,-h)$であり、その点曲線とtranslationは
\cite[Definition--Lemma 2.13]{WWY}の表示である。
第一点$p=(x,y,z)$に対する第二点の行列は
\[
 M(p)=\begin{pmatrix}-y&x&-hz\\-hx&-z&y\\z&-hy&-x\end{pmatrix},
 \qquad\det M=-h(x^3+y^3+z^3+h^2xyz).
\]
$E_h$の滑らかさからrankは常に2である。
rankが1以下なら全余因子が零で行列式の勾配も零となるからである。
第二射影も同様なので、双線形スキームは自己同型のgraphとなる。
$o$でのkernelが$[1:-1:-h]$であり、規約が確定する。
$h\ne0$によりtranslationは非恒等で、不動点を持たない。

中心座標が単元のchartでは中心交換関係により第二点は第一点と同じで、
残る式が$Y$を与える。
境界の$(p,\sigma_hp)$の近傍では、二列$p,p'$のある二次minorが単元である。
中心交換関係$up_i'-p_i u'=0$の二本は可逆な連立一次系となり$u=u'=0$を強制する。
$v=v'=0$も同様。この議論は冪零試験基底でも成立するので、余分な無限小方向はない。
従って双線形零スキームは$\Gamma_{\sigma_h}\sqcup\Delta_Y$となり、
graphを反復すれば点スキームとshiftを得る。

境界で二次関係を使うと
\[
 T=6xyz-3hx^3+3hy^3-3hz^3,\qquad
 3\Omega_h=(1+h^3)(x^3+z^3)+(1-h^3)y^3-2h^2xyz.
\]
\cite[(2.20)]{WWY}の標準中心元の特殊化は$g=-3h\Omega_h$であり、
同文献Lemma 2.19のTHCR表示から$\Omega_h$は普遍楕円点加群を零化する。
従って三次関係は楕円成分に追加条件を課さない。
定常成分では$w_f=0$となり、shiftは恒等なので全てのshifted条件が同じである。
\end{proof}
\begin{corollary}[埋込みとshiftを指定した点データによる復元]
指定した五座標、原点と正規形の下で、全点データは$(q_1,q_2,q_3,h,f)$を一意に回復する。
\end{corollary}
\begin{proof}
零次元成分と楕円成分は分離しており、前節の復元が$(q_1,q_2,q_3,h^2,f)$を与える。
残る符号は$\sigma_h(o)=[1:-1:-h]$から決まる。
$h\mapsto-h$で反対代数に移り、楕円shiftは逆向きになる。
抽象スキームだけ、座標を忘れた圏、標識なし導来圏の分類ではない。
\end{proof}
